\chapter{Implementacja systemu i środowiska pomiarowego}
\label{ch:implementacja_systemu}

W rozdziale tym szczegółowo opisano proces implementacji chmurowego środowiska badawczego. Przedstawiono strukturę repozytorium kodu, konfigurację manifestów deklaratywnych Kubernetes oraz zaawansowaną logikę skryptów testowych, uwzględniającą specyfikę bezpośredniego routingu sieciowego bez użycia serwerów DNS.

\section{Struktura repozytorium kodu źródłowego}
\label{sec:struktura_repozytorium}
Projekt badawczy został zorganizowany w jednolitą strukturę katalogów (ang. \textit{monorepo}), co zapewnia pełną powtarzalność oraz ułatwia automatyzację procesów budowania obrazów kontenerowych. Kluczowe elementy składowe projektu przedstawiają się następująco:

\begin{verbatim}
magisterka/
├── apps/
│   ├── cpu-service/       # Aplikacja CPU-bound (FastAPI + Pillow)
│   ├── io-service/        # Aplikacja I/O-bound (asyncio + httpx)
│   ├── echo-server/       # Serwis pomocniczy
│   └── memory-service/    # Rozszerzenie bazy badawczej (In-memory cache)
├── k8s/
│   ├── base/              # Definicja Namespace magisterka
│   ├── cpu-service/       # Manifesty Deployment, Service, HPA (x3)
│   ├── io-service/        # Manifesty Deployment, Service, HPA (x3)
│   └── monitoring/        # Konfiguracje Prometheus i Prometheus Adapter
└── tests/
    └── k6/                # Skrypty testowe (cpu-load-test.js, io-load-test.js)
\end{verbatim}

\section{Implementacja bezpośredniego routingu sieciowego w skryptach k6}
\label{sec:implementacja_routingu_k6}
W celu eliminacji dodatkowych warstw translacji nazw i uniknięcia narzutu konfiguracyjnego związanego z serwerami DNS w systemie operacyjnym generatora ruchu, mechanizm mapowania domen został wstrzyknięty bezpośrednio do kodu skryptów testowych \textit{k6}. 

Wykorzystano do tego celu zaawansowaną konfigurację pola \texttt{hosts} wewnątrz globalnego obiektu opcji środowiska uruchomieniowego \textit{k6}. Silnik pomiarowy automatycznie podmienia adresy IP na poziomie niskopoziomowych gniazd sieciowych, wstrzykując jednocześnie prawidłową wartość tekstową do nagłówka \texttt{Host} protokołu HTTP/1.1. Rozwiązanie to symuluje rzeczywiste zachowanie systemów produkcyjnych przechodzących przez chmurowy Load Balancer GCP. Pełną strukturę konfiguracyjną skryptu obciążeniowego dla serwisu obliczeniowego przedstawia listing \ref{lst:implementacja_k6_cpu}.

\begin{lstlisting}[language=JavaScript, caption={Konfiguracja bezpośredniego mapowania adresów IP Load Balancera GCP w skrypcie k6}, label={lst:implementacja_k6_cpu}]
import http from 'k6/http';
import { check, sleep } from 'k6';

export const options = {
  scenarios: {
    fibonacci_stress: {
      executor: 'ramping-vus',
      startVUs: 1,
      stages: [
        { duration: '60s', target: 50 },  // Faza Ramp-up
        { duration: '120s', target: 50 }, // Faza Steady-state
        { duration: '30s', target: 0 },   // Faza Ramp-down
      ],
      gracefulRampDown: '30s',
      exec: 'runFibonacci',
    },
  },
  hosts: {
    'cpu.saas-magisterka.internal': '35.242.210.14',
    'io.saas-magisterka.internal': '35.242.215.88',
  },
  thresholds: {
    'http_req_duration': ['p(95)<5000'], // Progi akceptacji SLA
    'http_req_failed': ['rate<0.01'],
  },
};

export function runFibonacci() {
  const url = 'http://cpu.saas-magisterka.internal/fibonacci?n=32';
  const res = http.get(url);
  check(res, {
    'status jest 200': (r) => r.status === 200,
  });
  sleep(0.5);
}
\end{lstlisting}

\section{Automatyzacja procesów telemetrycznych i eksport danych}
\label{sec:implementacja_telemetrii_grafana}
Akwizycja parametrów wydajnościowych została w pełni scentralizowana i zautomatyzowana w oparciu o zapytania API do bazy danych szeregów czasowych. Odstąpiono od parsowania logów terminalowych na rzecz ciągłego monitorowania parametrów infrastrukturalnych dostarczanych przez \texttt{kube-state-metrics} oraz \texttt{cAdvisor}. Interwał skrobania (ang. \textit{scrape interval}) został zaostrzony do wartości $\Delta t = 5$ sekund, co pozwala na uchwycenie gwałtownych anomalii chmurowych.

Dla każdego z 8 realizowanych eksperymentów, po zakończeniu pętli testowej, z platformy \textit{Grafana} za pomocą zautomatyzowanych zapytań HTTP pobierane były kompletne profile pomiarowe w formacie plików CSV. Zbiór ten obejmował pięć kluczowych wymiarów telemetrycznych:
\begin{enumerate}
    \item \textbf{Dynamiczna liczba replik HPA:} aktualna liczba Podów w stanie \texttt{Running} rejestrowana w dziedzinie czasu przez płaszczyznę kontrolną.
    \item \textbf{Zużycie CPU per Pod:} rzeczywiste obciążenie procesora wyrażone w milicząstkach rdzenia (\textit{mCPU}), mapowane bezpośrednio z liczników cgroups systemu Linux.
    \item \textbf{Alokacja pamięci RAM:} dynamiczne zmiany w zużyciu pamięci operacyjnej kontenerów wyrażone w megabajtach (MB).
    \item \textbf{Chwilowa przepustowość (RPS):} zagregowana liczba żądań HTTP obsłużonych pomyślnie przez punkty końcowe usług w każdej sekundzie testu.
    \item \textbf{Wykresy opóźnień sieciowych (Latency):} pełne przebiegi czasowe czasów odpowiedzi z podziałem na percentyle wydajnościowe $p_{50}$, $p_{95}$ oraz $p_{99}$.
\end{enumerate}

\section{Wdrożenie manifestów orkiestracji kontenerowej}
\label{sec:implementacja_manifestow_k8s}
Wszystkie zasoby klastra GKE są definiowane w sposób deklaratywny. Przykładową implementację obiektu typu \texttt{Deployment} dla serwisu obliczeniowego, zawierającą ścisłe limity zasobowe cgroups (\texttt{resources.limits}) oraz konfiguracje sond gotowości sieciowej, przedstawia listing \ref{lst:k8s_deployment_cpu}.

\begin{lstlisting}[language=XML, caption={Manifest deklaratywny Deployment dla cpu-service na platformie GKE}, label={lst:k8s_deployment_cpu}]
apiVersion: apps/v1
kind: Deployment
metadata:
  name: cpu-service
  namespace: magisterka
  labels:
    app: cpu-service
spec:
  replicas: 2
  selector:
    matchLabels:
      app: cpu-service
  template:
    metadata:
      labels:
        app: cpu-service
    spec:
      containers:
      - name: cpu-service
        image: europe-west3-docker.pkg.dev/saas-project/gar/cpu-service:v1.0
        imagePullPolicy: IfNotPresent
        ports:
        - containerPort: 8000
        resources:
          requests:
            cpu: "250m"
            memory: "256Mi"
          limits:
            cpu: "1000m"
            memory: "512Mi"
        livenessProbe:
          httpGet:
            path: /health
            port: 8000
          initialDelaySeconds: 5
          periodSeconds: 10
        readinessProbe:
          httpGet:
            path: /health
            port: 8000
          initialDelaySeconds: 5
          periodSeconds: 5
\end{lstlisting}